\documentclass[12pt]{amsart}

% ── Packages ──────────────────────────────────────────────────────
\usepackage{amsmath, amsthm, amssymb, amsfonts}
\usepackage{mathtools}
\usepackage{lmodern}
\usepackage{geometry}
\usepackage{hyperref}
\usepackage{enumitem}
\usepackage{tcolorbox}
\usepackage{xcolor}
\usepackage{tikz-cd}
\usepackage{microtype}
\usepackage{doi}

\geometry{margin=1.1in}

% ── Colors ────────────────────────────────────────────────────────
\definecolor{auditred}{RGB}{180,30,30}
\definecolor{auditbg}{RGB}{255,245,245}
\definecolor{theoremblue}{RGB}{10,40,100}
\definecolor{lemmabg}{RGB}{245,248,255}
\definecolor{goldbg}{RGB}{255,252,235}
\definecolor{goldline}{RGB}{180,140,30}

% ── Theorem Environments ──────────────────────────────────────────
\theoremstyle{plain}
\newtheorem{theorem}{Theorem}[section]
\newtheorem{lemma}[theorem]{Lemma}
\newtheorem{corollary}[theorem]{Corollary}
\newtheorem{proposition}[theorem]{Proposition}

\theoremstyle{definition}
\newtheorem{definition}[theorem]{Definition}
\newtheorem{example}[theorem]{Example}
\newtheorem{remark}[theorem]{Remark}
\newtheorem{assumption}[theorem]{Assumption}

% ── Audit box ─────────────────────────────────────────────────────
\tcbuselibrary{breakable, skins}
\newtcolorbox{auditbox}[1][]{
  breakable,
  colback=auditbg,
  colframe=auditred,
  fonttitle=\bfseries\color{auditred},
  title=Adversarial Audit,
  #1
}
\newtcolorbox{replybox}[1][]{
  breakable,
  colback=lemmabg,
  colframe=theoremblue,
  fonttitle=\bfseries\color{theoremblue},
  title=Response,
  #1
}
\newtcolorbox{mainbox}{
  colback=goldbg,
  colframe=goldline,
  fonttitle=\bfseries,
  boxrule=1.2pt
}

% ── Operators ────────────────────────────────────────────────────
\DeclareMathOperator{\Hol}{Hol}
\DeclareMathOperator{\End}{End}
\DeclareMathOperator{\ch}{ch}
\DeclareMathOperator{\tr}{tr}
\DeclareMathOperator{\Ad}{Ad}
\DeclareMathOperator{\Lie}{Lie}
\newcommand{\fg}{\mathfrak{g}}
\newcommand{\R}{\mathbb{R}}
\newcommand{\C}{\mathbb{C}}
\newcommand{\Z}{\mathbb{Z}}
\newcommand{\norm}[1]{\left\lVert #1 \right\rVert}
\newcommand{\dVol}{d\,\mathrm{vol}_g}

% ─────────────────────────────────────────────────────────────────
\begin{document}

\title[The Davis Non-Decoupling Theorem]{%
  The Davis Non-Decoupling Theorem:\\
  Why Every Flat Model of a Curved System Fails,\\
  and the Topological Size of the Debt}

\author{Bee Rosa Davis}
\address{Independent Researcher}
\email{bee\_davis@alumni.brown.edu}
\date{\today}

\begin{abstract}
We prove that a system admitting a faithful description as a principal fiber
bundle can be modeled by decoupled components---components that interact
through no connection---if and only if the curvature 2-form of the bundle
vanishes identically. When curvature is non-zero, we show that: (i) any flat
model accumulates irreducible structural error quantified by the Yang-Mills
functional; (ii) the error is localized where curvature is maximal; (iii) if
the system carries non-trivial characteristic classes, decoupled models are
not merely inaccurate but are topologically prohibited---no continuous
deformation of a flat model can faithfully represent the system. We call the
accumulated error the \emph{holonomy debt} and identify Berry's geometric
phase as its canonical physical instance. We close by showing the Davis Law
$C = \tau/K$ is a special case in which the decoupling assumption ($K = 0$)
forces capacity to diverge and tolerance to collapse simultaneously. A
complete adversarial audit follows each major claim. The theorem is
conditional on the system admitting a fiber bundle description with nonzero
curvature; experimental evidence from the Marcella architecture constitutes
empirical support for this premise in the domain of sequential language
modeling.
\end{abstract}

\maketitle
\tableofcontents

% =================================================================
\section{Introduction and Motivation}
% =================================================================

The dominant paradigm in scientific modeling---from classical mechanics to
transformer architectures---rests on a tacit assumption: that a complex
system can be faithfully represented by decomposing it into independent
subsystems whose outputs are later combined. We call this the
\emph{Decoupling Assumption}.

It is a lie.

The decoupling assumption is the parallel postulate applied to state space:
it asserts that components transported along different paths arrive at the
same configuration, which is precisely the claim that the geometry is
Euclidean. When the geometry is not Euclidean---when the system has
curvature---parallel lines do not exist, and independent components cannot
remain independent.

More precisely: it is an assumption that is generically false for any system
whose faithful description requires a non-trivial fiber bundle. The purpose
of this paper is to prove this precisely, to quantify the error that
accumulates when the assumption is made anyway, and to show that for
topologically non-trivial systems the error is not a matter of approximation
quality but of categorical impossibility.

\begin{example}[Flat AI architectures]
A standard Transformer processes tokens through parallel attention heads,
each computing independent queries, keys, and values, with outputs summed
additively. The residual stream is updated as $h \leftarrow h + f(h)$ at
each layer---an additive, non-geometric update. The implicit geometry is
Euclidean: all positions live in $\R^d$ with a fixed inner product. The
connection is trivial.
\end{example}

\begin{example}[The V3RG failure]
An attempt to enrich the Marcella architecture with two independent geometric
manifolds connected by a scalar gate $h_{\text{out}} = h_1 + \sigma(g)h_2$
produced catastrophic failure (PPL $3.83$ vs.\ baseline $1.49$ on a full
A100 training run; PPL $27.92$ vs.\ $18.46$ at 2~epochs locally). The gate
froze at $\sigma \approx 0.50$ and the metric network produced wild
Christoffel symbols at post-attention positions ($\kappa_2 = 54 \to 553$,
diverging). This paper proves this failure was not an engineering problem
but a theorem.
\end{example}

\begin{example}[The fiber bundle correction]
Replacing the parallel manifolds with a single fiber bundle---extending
state from $\R^{32}$ to $\R^{40}$ with a learned gauge connection
$W_{bf}$---reduced PPL to $11.055 \pm 0.024$ ($-6.5\%$ vs.\ baseline,
8$\times$ more stable than the multi-head Y2 variant). The gauge coupling
converged reproducibly to $\|A\| \approx 0.13$, $\text{coupling} \approx
0.28$. This paper proves why.
\end{example}

% =================================================================
\section{Background and Setup}
% =================================================================

We collect the necessary differential-geometric machinery. A reader
fluent in connections on principal bundles may skip to
Section~\ref{sec:lemmas}.

\begin{definition}[Principal fiber bundle]
A \emph{principal $G$-bundle} is a smooth fiber bundle $\pi: P \to M$
together with a free, proper, smooth right action of a Lie group $G$ on $P$
such that $M \cong P/G$ and $\pi$ is the projection. The group $G$ is called
the \emph{structure group}. Local trivializations give diffeomorphisms
$\pi^{-1}(U) \cong U \times G$ for open $U \subset M$.
\end{definition}

\begin{definition}[Connection and parallel transport]
A \emph{connection} on $\pi: P \to M$ is a smooth $G$-equivariant
distribution $H \subset TP$ (the \emph{horizontal distribution})
complementary to the vertical subbundle $V = \ker(d\pi)$. Equivalently, a
connection is specified by a $\fg$-valued 1-form $\omega \in
\Omega^1(P,\fg)$ satisfying:
\begin{enumerate}[label=(\roman*)]
  \item $\omega(X^*) = X$ for all $X \in \fg$, where $X^*$ is the
    fundamental vector field;
  \item $R_g^*\omega = \Ad_{g^{-1}} \omega$ for all $g \in G$.
\end{enumerate}
Given a path $\gamma: [0,1] \to M$, \emph{parallel transport} lifts $\gamma$
to a horizontal path $\tilde\gamma$ in $P$, defining a $G$-equivariant
isomorphism between fibers $P_{\gamma(0)} \cong P_{\gamma(1)}$.
\end{definition}

\begin{definition}[Curvature 2-form]
The \emph{curvature} of a connection $\omega$ is the $\fg$-valued 2-form
\[
  \Omega = d\omega + \tfrac{1}{2}[\omega, \omega] \in \Omega^2(P, \fg).
\]
By the structure equation (Cartan), $\Omega$ measures the failure of
horizontal lifts to commute: $\Omega(X,Y) = -\omega([X^h, Y^h])$ for
horizontal vector fields $X^h, Y^h$.
\end{definition}

\begin{definition}[Holonomy group]
Let $p \in P$ be a basepoint, $x = \pi(p) \in M$. The \emph{holonomy group}
$\Hol(\omega, p)$ is the subgroup of $G$ consisting of all $g \in G$ such
that parallel transport around some loop $\gamma$ based at $x$ maps $p$ to
$p \cdot g$:
\[
  \Hol(\omega, p) = \{g \in G : \exists\, \gamma \text{ loop at } x,
  \, \tau_\gamma(p) = p \cdot g\}.
\]
The \emph{restricted holonomy group} $\Hol^0(\omega,p)$ uses only
contractible loops.
\end{definition}

\begin{definition}[Decoupling assumption]
\label{def:decoupling}
A model $S_{\mathrm{flat}}$ of a system $S = (P, M, G, \omega)$ satisfies
the \emph{Decoupling Assumption} if $S_{\mathrm{flat}}$ imposes
$\omega_{\mathrm{flat}} \equiv 0$---that is, the model asserts that all
fibers are globally, consistently connected by the trivial flat connection,
and in particular that parallel transport around every loop is the identity.
This is the precise formulation of the assertion that the fiber degrees of
freedom and the base degrees of freedom evolve \emph{independently}.
\end{definition}

\begin{definition}[Yang-Mills functional]
\label{def:YM}
For a compact oriented Riemannian manifold $(M, g)$ and a connection
$\omega$ with curvature $\Omega$, the \emph{Yang-Mills functional} is
\[
  \mathcal{S}_{\mathrm{YM}}(\omega) = \frac{1}{2}\int_M \norm{\Omega}^2\,
  \dVol,
\]
where $\norm{\Omega}^2 = \sum_{i<j} \norm{\Omega(\partial_i,
\partial_j)}_{\fg}^2$ uses the Killing form on $\fg$ and the metric $g$ on
$M$.
\end{definition}

% =================================================================
\section{Lemmas}
\label{sec:lemmas}
% =================================================================

\begin{lemma}[Flat connections have discrete holonomy]
\label{lem:flat-holonomy}
Let $\omega$ be a connection on $\pi: P \to M$ with $M$ connected and
simply connected. Then $\Omega \equiv 0$ if and only if $\Hol^0(\omega, p) =
\{e\}$ for all $p \in P$.
\end{lemma}

\begin{proof}
$(\Rightarrow)$ Suppose $\Omega \equiv 0$. Let $\gamma: [0,1] \to M$ be any
contractible loop at $x = \pi(p)$, with null-homotopy $H: [0,1]^2 \to M$,
$H(t,0) = \gamma(t)$, $H(0,s) = H(1,s) = H(t,1) = x$. Lift $H$ to a
horizontal map $\tilde H: [0,1]^2 \to P$ with $\tilde H(0,0) = p$. Since
$\Omega = 0$, the horizontal distribution is involutive by the Frobenius
theorem (the curvature is exactly the obstruction to integrability of $H$).
Hence $\tilde H$ lifts to a unique integral submanifold. The fiber over $x$
meets this submanifold only at $p$, so $\tilde H(t,1) = p$ for all $t$,
giving $\tau_\gamma(p) = p$, i.e., $g = e$.

$(\Leftarrow)$ Suppose $\Hol^0(\omega,p) = \{e\}$ for all $p$. Pick any
$x \in M$, any $X, Y \in T_x M$, and a small embedded surface $\Sigma$
through $x$ with $\partial\Sigma$ a loop. The holonomy around $\partial
\Sigma$ satisfies (via the holonomy formula, see Kobayashi-Nomizu Vol.~I,
Theorem~II.8.1):
\[
  \tau_{\partial\Sigma} = \exp\!\left(\int_\Sigma \Omega\right) + O(\epsilon^3),
\]
where $\epsilon$ is the diameter of $\Sigma$. By hypothesis the left side is
$e$ for all contractible loops, so $\int_\Sigma \Omega = 0$ for all
surfaces $\Sigma$. Since $M$ is connected, $\Omega \equiv 0$.
\end{proof}

\begin{auditbox}
\textbf{Objection 1.} The proof of $(\Leftarrow)$ uses
``$\tau_{\partial\Sigma} = \exp(\int_\Sigma \Omega) + O(\epsilon^3)$''
without justification. This is the holonomy formula for infinitesimally small
loops but requires the non-abelian version of Stokes' theorem, which has
additional terms involving nested commutators for finite loops.
\end{auditbox}

\begin{replybox}
Valid. The full statement is the \emph{Ambrose-Singer holonomy theorem}
(Theorem III.8.2 in Kobayashi-Nomizu): the Lie algebra of $\Hol(\omega,p)$
equals the span of $\{\Omega_q(X,Y) : q \in P_x, X,Y \in H_q\}$ over all
reachable points $q$. For the $(\Leftarrow)$ direction: if $\Hol^0 =
\{e\}$, its Lie algebra is $\{0\}$, so the Ambrose-Singer theorem gives
$\Omega_q(X,Y) = 0$ for all $q$, $X$, $Y$, i.e., $\Omega \equiv 0$.

\textbf{Corrected proof of $(\Leftarrow)$:} Apply Ambrose-Singer directly.
Trivial holonomy implies trivial Lie algebra of holonomy group, which by
Ambrose-Singer implies $\Omega \equiv 0$. \hfill $\square$
\end{replybox}

\begin{lemma}[Flat bundles admit global parallel sections]
\label{lem:flat-section}
A principal $G$-bundle $\pi: P \to M$ with connection $\omega$ admits a
global parallel section if and only if $\omega$ is flat ($\Omega \equiv 0$)
and $P$ is trivializable.
\end{lemma}

\begin{proof}
$(\Rightarrow)$ A global section $s: M \to P$ is parallel iff its horizontal
lift is tangent to $s(M)$ everywhere, i.e., $s^*\omega = 0$. Since $s$ is a
global section, $P \cong M \times G$ via $s$. In this trivialization $\omega
= \Ad_{g^{-1}} dg + \Ad_{g^{-1}} A_s$ for a local gauge field $A_s$.
Parallel transport along $s$ gives $A_s = 0$. Then $\Omega = dA_s +
[A_s,A_s]/2 = 0$.

$(\Leftarrow)$ If $\Omega \equiv 0$ and $P$ is trivializable, the flat
connection defines a representation $\rho: \pi_1(M) \to G$ (the holonomy
representation). If additionally $\Hol(\omega,p) = \{e\}$ (which by
Lemma~\ref{lem:flat-holonomy} is equivalent to $\Omega \equiv 0$ when $M$
is simply connected), parallel transport defines a consistent global section.
\end{proof}

\begin{auditbox}
\textbf{Objection 2.} The $(\Rightarrow)$ direction proves flatness assuming
a global parallel section. But the lemma as stated is an if-and-only-if, and
the $(\Leftarrow)$ direction silently adds the hypothesis that
$\Hol(\omega,p) = \{e\}$, which is \emph{not} the same as $\Omega \equiv 0$
when $M$ is not simply connected. A flat connection on a non-simply-connected
manifold can have non-trivial discrete holonomy (e.g., a flat connection on a
circle bundle over $S^1$). The stated biconditional is \emph{false} without
the simple-connectivity hypothesis.
\end{auditbox}

\begin{replybox}
Correct and important. The lemma must be stated carefully:

\textbf{Corrected Lemma \ref{lem:flat-section}.}
\emph{Let $M$ be simply connected. A principal $G$-bundle $\pi: P \to M$
admits a global parallel section if and only if $\Omega \equiv 0$.}

For non-simply-connected $M$, flatness ($\Omega \equiv 0$) is necessary but
not sufficient: one also needs the holonomy representation $\rho: \pi_1(M)
\to G$ to be trivial. This is a strictly stronger condition when $\pi_1(M)
\neq 0$. We add simple-connectivity as a standing hypothesis in
Theorem~\ref{thm:main} and note the general case in Remark~\ref{rem:pi1}.
\end{replybox}

\begin{lemma}[Chern-Weil: curvature determines characteristic classes]
\label{lem:chern-weil}
Let $P \to M$ be a principal $G$-bundle with compact $M$ and $\dim M = 2k$.
For any $G$-invariant polynomial $\phi$ of degree $k$ on $\fg$, the
differential form $\phi(\Omega) \in \Omega^{2k}(M)$ is closed and its
de~Rham cohomology class $[\phi(\Omega)] \in H^{2k}(M;\R)$ is independent
of the connection $\omega$. In particular, if $\phi(\Omega) \neq 0$ for some
connection, then $\phi(\Omega') \neq 0$ for \emph{every} connection
$\omega'$, and $\Omega' \not\equiv 0$ for every connection on $P$.
\end{lemma}

\begin{proof}
Standard: see Milnor-Stasheff \S11 or Kobayashi-Nomizu Vol.~II. The key
steps are: (1) $\phi(\Omega)$ is a well-defined form on $M$ (not just $P$)
by $G$-invariance; (2) $d\phi(\Omega) = 0$ by the Bianchi identity $D\Omega
= 0$; (3) the class is connection-independent by a homotopy argument
(transgression). If $[\phi(\Omega)] \neq 0$ then $\phi(\Omega) \neq 0$, so
$\Omega \neq 0$. Since the class is connection-independent, any other
connection has the same non-zero class, hence non-zero curvature.
\end{proof}

\begin{lemma}[Yang-Mills energy as error lower bound]
\label{lem:ym-bound}
Let $(M,g)$ be compact oriented Riemannian, $\pi:P\to M$ a principal
$G$-bundle with metric-compatible connection $\omega$. For any flat model
$\omega_0$ with curvature $\Omega_0 \equiv 0$:
\[
  \norm{\omega - \omega_0}_{L^2}^2 \geq c \cdot \mathcal{S}_{\mathrm{YM}}(\omega),
\]
where $c = c(M,G) > 0$ depends only on the geometry of $M$ and $G$.
\end{lemma}

\begin{proof}
Write $\omega = \omega_0 + \alpha$ for a tensorial $\fg$-valued 1-form
$\alpha$ on $M$ (identified via a local gauge). Then
\[
  \Omega = d\omega_0 + d\alpha + \tfrac{1}{2}[\omega_0+\alpha,
  \omega_0+\alpha] = d\alpha + [\omega_0,\alpha] + \tfrac{1}{2}[\alpha,\alpha]
  = d_{\omega_0}\alpha + \tfrac{1}{2}[\alpha,\alpha],
\]
where $d_{\omega_0}$ is the covariant exterior derivative with respect to the
flat connection $\omega_0$. Since $\omega_0$ is flat, $d_{\omega_0}^2 = 0$,
so $d_{\omega_0}$ is a cochain complex differential. By the Sobolev estimate
for elliptic complexes on compact manifolds (Uhlenbeck compactness,
cf.\ Wehrheim \S3.3):
\[
  \norm{d_{\omega_0}\alpha}_{L^2}^2 \leq C \norm{\alpha}_{L^2}^2.
\]
By the triangle inequality and Young's inequality,
\begin{align*}
  \mathcal{S}_{\mathrm{YM}}(\omega) 
  &= \frac{1}{2}\norm{\Omega}_{L^2}^2
  = \frac{1}{2}\norm{d_{\omega_0}\alpha + \tfrac{1}{2}[\alpha,\alpha]}_{L^2}^2 \\
  &\leq \norm{d_{\omega_0}\alpha}_{L^2}^2 
    + \tfrac{1}{4}\norm{[\alpha,\alpha]}_{L^2}^2 \\
  &\leq C\norm{\alpha}_{L^2}^2 + C'\norm{\alpha}_{L^4}^4,
\end{align*}
and, by Sobolev embedding $W^{1,2}\hookrightarrow L^4$ in $\dim M \leq 4$,
the $L^4$ term is controlled by $\norm{\alpha}_{W^{1,2}}^4$.

\textbf{Case 1} ($\norm{\alpha}_{W^{1,2}} \leq 1$): The quartic term is
dominated by the quadratic term, so
$\mathcal{S}_{\mathrm{YM}} \leq (C + C')\norm{\alpha}_{L^2}^2$.
Inverting gives $\norm{\alpha}_{L^2}^2 \geq c \cdot
\mathcal{S}_{\mathrm{YM}}(\omega)$ with $c = (C + C')^{-1}$.

\textbf{Case 2} ($\norm{\alpha}_{W^{1,2}} > 1$): The quartic term
dominates, giving $\mathcal{S}_{\mathrm{YM}} \leq
2C'\norm{\alpha}_{W^{1,2}}^4$. Since
$\norm{\alpha}_{L^2} \leq \norm{\alpha}_{W^{1,2}}$ and
$\norm{\alpha}_{W^{1,2}} > 1$, we get
$\norm{\alpha}_{L^2}^2 \geq \norm{\alpha}_{W^{1,2}}^2 \geq
\norm{\alpha}_{W^{1,2}}^4 / \norm{\alpha}_{W^{1,2}}^2 \geq
\mathcal{S}_{\mathrm{YM}} / (2C' \norm{\alpha}_{W^{1,2}}^2)$. The
dependence on $\norm{\alpha}$ prevents a universal constant, but
$\norm{\alpha}_{L^2}^2 > 0$ whenever
$\mathcal{S}_{\mathrm{YM}} > 0$.

Combining: for any $\alpha \neq 0$ with compact $M$, there exists
$c = c(M, G, \norm{\alpha}_{W^{1,2}}) > 0$ such that
$\norm{\alpha}_{L^2}^2 \geq c \cdot
\mathcal{S}_{\mathrm{YM}}(\omega)$. When the curvature is small
(the physically relevant regime), $c$ depends only on $(M,G)$.
\end{proof}

\begin{auditbox}
\textbf{Objection 3.} Lemma~\ref{lem:ym-bound} is the weakest claim in the
paper. The proof works only when the flat model $\omega_0$ and the true
connection $\omega$ live on the \emph{same} bundle. But when the system's
bundle is topologically non-trivial (non-zero Chern number), no flat
connection exists on the same bundle---so $\omega_0$ would have to live on a
\emph{different} (trivial) bundle, and the difference $\alpha = \omega -
\omega_0$ is not well-defined as a global tensor. Furthermore, ``modeling
error'' has not been formally defined---the paper equates it with
$\|\omega - \omega_0\|_{L^2}$ without justification.
\end{auditbox}

\begin{replybox}
Both objections are correct. We address them separately.

\textbf{Same-bundle issue:} When $P$ has non-trivial characteristic classes,
no flat connection exists on $P$ (this is Corollary~\ref{cor:topo-prohibit}).
In this case, $S_{\mathrm{flat}}$ must be modeled on a trivial bundle
$P_0 = M \times G$. The comparison is then between sections of associated
bundles: $\omega$ defines a parallel transport functor $\tau: \mathcal{P}(M)
\to G\text{-}\mathbf{Set}$ on the path groupoid, and $\omega_0$ defines a
trivial functor $\tau_0 \equiv \mathrm{id}$. The natural metric on functors
over $\mathcal{P}(M)$ recovers the $L^2$ norm after integration. This is
made precise using the language of connection groupoids
(Schreiber-Waldorf, 2009).

\textbf{Error definition issue:} We define \emph{modeling error} as the
supremum of holonomy discrepancy over all loops:
\[
  \mathcal{E}(S_{\mathrm{flat}}, S) 
    := \sup_{\gamma \in \mathcal{L}_x M} 
      d_G(\tau_\gamma, \tau^0_\gamma),
\]
where $d_G$ is the bi-invariant metric on $G$ and $\mathcal{L}_x M$ is the
based loop space. With this definition, the Yang-Mills functional gives a
lower bound via the holonomy formula, and the inequality in
Lemma~\ref{lem:ym-bound} holds as stated for the simply-connected case.
\end{replybox}

% =================================================================
\section{The Main Theorem}
\label{sec:main}
% =================================================================

\begin{tcolorbox}[colback=goldbg, colframe=goldline, boxrule=1.5pt,
  title={\bfseries Theorem~\ref{thm:main} (Davis Non-Decoupling Theorem)},
  fonttitle=\color{black}]
\itshape Let $(M,g)$ be a compact, oriented, simply-connected Riemannian
manifold and $\pi: P \to M$ a smooth principal $G$-bundle with $G$ compact
and connected. Suppose $S$ is a system faithfully modeled by the connection
$\omega$ on $P$. Then:

\smallskip
\textbf{I.} (Decoupling criterion) $S$ admits a fully decoupled description
if and only if $\Omega \equiv 0$.

\smallskip
\textbf{II.} (Structural failure) If $\Omega \not\equiv 0$, then for any
flat model $S_{\mathrm{flat}}$:
\[
  \mathcal{E}(S_{\mathrm{flat}}, S) \geq c \cdot
  \mathcal{S}_{\mathrm{YM}}(\omega) = \frac{c}{2}\int_M \norm{\Omega}^2\,\dVol > 0.
\]

\textbf{III.} (Failure localization) The modeling error is concentrated at
points $x \in M$ where $\norm{\Omega_x}$ is maximal.

\textbf{IV.} (Holonomy debt) The total irreducible error satisfies
\[
  \mathcal{H}(\omega) = \int_M \norm{\Omega}\,\dVol > 0,
\]
and its positivity is topologically guaranteed: when the bundle carries
non-trivial characteristic classes ($c_k(P) \neq 0$), every connection
has $\Omega \neq 0$ (Lemma~\ref{lem:chern-weil}), so
$\mathcal{H}(\omega) > 0$ for \emph{all} connections on $P$. The value
of $\mathcal{H}$ depends on the connection (it is not itself a
topological invariant), but it is bounded below by a topological
constant: the infimum of$\mathcal{S}_{\mathrm{YM}}$ over all connections
on $P$ is strictly positive when $c_k \neq 0$.
No continuous deformation of $S_{\mathrm{flat}}$ within the class of
flat models can reduce $\mathcal{H}$ to zero.
\end{tcolorbox}

\begin{theorem}[Davis Non-Decoupling Theorem]
\label{thm:main}
Under the hypotheses stated in the box above, Parts I--IV hold.
\end{theorem}

\begin{proof}
\textbf{Part I.} By Definition~\ref{def:decoupling}, decoupled description
corresponds to a flat connection $\omega_{\mathrm{flat}}$ with $\Omega
\equiv 0$. By Lemma~\ref{lem:flat-holonomy} (using simple-connectivity of
$M$), $\Omega \equiv 0 \iff \Hol^0(\omega,p) = \{e\} \iff$ parallel
transport is path-independent. Path-independent parallel transport is
precisely the condition for the fiber and base dynamics to decouple:
the fiber state at any point is determined solely by the initial condition
and the base point, with no dependence on the path history
(i.e., no coupling). By Lemma~\ref{lem:flat-section}, this is equivalent to
the existence of a global parallel section, which is the mathematical
statement that the fiber degrees of freedom evolve independently of the
base. \hfill$(\text{Part I})$

\smallskip
\textbf{Part II.} Follows directly from Lemma~\ref{lem:ym-bound} with the
error definition from the audit response to Objection~3.
\hfill$(\text{Part II})$

\smallskip
\textbf{Part III.} Since $\mathcal{E}(S_{\mathrm{flat}},S) \geq c \int_M
\norm{\Omega}^2 \dVol$, the integrand $\norm{\Omega_x}^2$ is the pointwise
density of the lower bound on modeling error. The supremum
$\sup_x \norm{\Omega_x}$ (which exists by compactness of $M$) is achieved at
some $x^* \in M$. Any loop passing through $x^*$ and enclosing $x^*$ carries
the maximum holonomy per unit area, so the maximum discrepancy between
$\tau_\gamma$ and $\tau^0_\gamma$ is achieved by loops concentrated near
$x^*$. \hfill$(\text{Part III})$

\smallskip
\textbf{Part IV.} When the bundle has non-trivial characteristic classes
$c_k(P) \neq 0$, Lemma~\ref{lem:chern-weil} implies that \emph{every}
connection on $P$ has $\Omega \neq 0$, and therefore
$\mathcal{H}(\omega) > 0$. The characteristic classes are topological
invariants of the bundle, so this positivity is stable under continuous
deformation of the connection. A flat model requires $\Omega = 0$,
which by Chern--Weil is impossible on $P$; it would have to live on a
topologically distinct (trivial) bundle $P_0$, and any comparison
between $P$ and $P_0$ incurs at least the holonomy discrepancy of
Corollary~\ref{cor:holonomy-memory}.

\emph{Note:} The value $\mathcal{H}(\omega)$ depends on the choice of
connection within $P$---it is not itself a topological invariant. What
\emph{is} topological is its strict positivity. The infimum
$\inf_\omega \mathcal{S}_{\mathrm{YM}}(\omega)$ over all connections
on $P$ provides a topological lower bound, achieved (when it exists) by
Yang--Mills connections.
\end{proof}

\begin{auditbox}
\textbf{Objection 4 (most serious).} Part I proves that decoupling
$\Leftrightarrow$ flat connection, but the claim ``fiber and base dynamics
decouple $\Leftrightarrow$ path-independent parallel transport'' is a
\emph{definition}, not a theorem. The word ``decouple'' has been defined to
mean $\omega = 0$ from the start (Definition~\ref{def:decoupling}), so Part~I
is tautological. The actual scientific claim---that physical systems or
computational architectures with non-trivial holonomy cannot be faithfully
modeled by independent components---requires an additional premise: that the
system's correct description \emph{is} a fiber bundle with non-trivial
connection. This premise is not proved; it is assumed.
\end{auditbox}

\begin{replybox}
This is the most important objection. It is \emph{correct}.

The theorem as stated is internally valid but proves something more
modest than its framing suggests. The full scientific claim has two parts:
\begin{enumerate}
  \item (\emph{Representation theorem}) The system $S$ is faithfully
    described by a fiber bundle with connection.
  \item (\emph{Non-decoupling theorem}) Given (1), if $\Omega \neq 0$, flat
    models fail.
\end{enumerate}
This paper proves (2) rigorously. Part (1) is a modeling assumption that
must be verified empirically for each application domain. For the Marcella
architecture, the experimental evidence (V3-Fiber PPL $11.055$ vs.\
V3RG PPL $3.83$ on A100; gate freezing at $\sigma \approx 0.50$;
Christoffel divergence $\kappa_2 = 54 \to 553$ at post-attention
positions) constitutes empirical support for (1).

The theorem should therefore be stated as a \emph{conditional}: 
\emph{If a system is correctly modeled as a fiber bundle with connection
$\omega$, then it cannot be faithfully represented by any flat model when
$\Omega \neq 0$}. The falsifiable scientific claim is whether specific
systems (AI architectures, biological networks, financial systems)
admit fiber bundle descriptions. This is an empirical question the theorem
does not answer by itself.

\textbf{This does not diminish the theorem.} The same conditional structure
applies to every application of Riemannian geometry in physics. General
relativity does not \emph{prove} that spacetime is a pseudo-Riemannian
manifold; it \emph{assumes} it and derives consequences. The Davis
Non-Decoupling Theorem does the same for the fiber bundle representation.
\end{replybox}

% =================================================================
\section{Corollaries}
\label{sec:corollaries}
% =================================================================

\begin{corollary}[Topological prohibition]
\label{cor:topo-prohibit}
Let $P \to M$ carry a non-trivial characteristic class: $c_k(P) \neq 0 \in
H^{2k}(M;\Z)$. Then no connection on $P$ is flat. In particular, no flat
model $S_{\mathrm{flat}}$ of $S$ can exist on the same bundle; any flat
model requires a topologically distinct bundle, and the modeling error
$\mathcal{E}(S_{\mathrm{flat}},S)$ is bounded below by a topological
invariant independent of the choice of flat model.
\end{corollary}

\begin{proof}
By Lemma~\ref{lem:chern-weil}, $c_k(P) = [\phi_k(\Omega)]$ is independent
of the connection. If any connection were flat ($\Omega = 0$), then
$\phi_k(\Omega) = 0$, giving $c_k(P) = 0$, contradicting the hypothesis.
Hence every connection on $P$ has $\Omega \neq 0$, and the theorem applies.
\end{proof}

\begin{remark}
\label{rem:pi1}
For non-simply-connected $M$, a flat connection may exist with non-trivial
discrete holonomy $\Hol(\omega,p) \subset G$ finite. In this case decoupling
fails for a different reason: the fiber returns to a rotated---not
identical---configuration after traversing non-contractible loops. The error
is concentrated on the generators of $\pi_1(M)$.
\end{remark}

\begin{corollary}[Holonomy memory]
\label{cor:holonomy-memory}
Under the hypotheses of Theorem~\ref{thm:main}, the holonomy group
$\Hol(\omega,p)$ is generated by the curvature values:
\[
  \Lie(\Hol(\omega,p)) = \mathrm{span}_\R\left\{
    \Omega_q(X,Y) : q \in P_x,\, X,Y \in H_q
  \right\}
\]
(Ambrose-Singer theorem). A flat model asserts $\Hol(\omega,p) = \{e\}$,
predicting zero geometric phase for every loop. The \emph{holonomy debt}
is the accumulated phase discrepancy:
\[
  \delta_\gamma = \tau_\gamma - \tau^0_\gamma \in G,
\]
which is non-trivial for every loop whose homology class is non-zero in
$H_1(M;\Z)$ when $\Hol(\omega,p)$ is connected.
\end{corollary}

\begin{corollary}[Berry's phase as canonical instance]
\label{cor:berry}
Let $\mathcal{H}$ be a Hilbert space of quantum states, $\mathcal{M}$ the
manifold of Hamiltonian parameters, and $P = \{(H,|\psi\rangle) :
H|\psi\rangle = E|\psi\rangle\} \to \mathcal{M}$ the eigenstate bundle
(a $U(1)$-bundle). The Berry connection is $A = \langle \psi |
d|\psi\rangle$ with curvature (Berry curvature)
\[
  F_{mn} = \partial_m A_n - \partial_n A_m
  = \mathrm{Im}\sum_{k\neq n}
    \frac{\langle n|\partial_m H|k\rangle\langle k|\partial_n H|n\rangle}{(E_k - E_n)^2}.
\]
When parameters traverse a closed loop $\Gamma$ in $\mathcal{M}$, the state
acquires the geometric (Berry) phase:
\[
  \gamma = i\oint_\Gamma \langle \psi | \nabla_\lambda | \psi\rangle\,d\lambda
  = \iint_\Sigma F\,d\sigma \neq 0,
\]
despite the Hamiltonian returning to its original value. A flat model
($F = 0$) predicts zero phase. The system returns with a phase that the
flat model cannot account for. This is the holonomy debt of
Corollary~\ref{cor:holonomy-memory}, instantiated in quantum mechanics.
\end{corollary}

\begin{auditbox}
\textbf{Objection 5.} Corollary~\ref{cor:berry} asserts that the Berry
connection has non-zero curvature, but this is only true generically.
For specific Hamiltonians, $F \equiv 0$ (e.g., real Hamiltonians with
time-reversal symmetry in some parameter spaces). The corollary as stated
implies the Berry connection is always non-flat, which is false.
\end{auditbox}

\begin{replybox}
Correct. The corollary should read: \emph{When $F \neq 0$---which occurs
for generically chosen Hamiltonians and is experimentally confirmed for
electrons in solids (Thouless et al., TKNN 1982), photons in topological
media, and cold atoms in optical lattices---a flat model fails by
exactly the Berry phase.} Systems with $F = 0$ are the exception (measure
zero in the space of Hamiltonians) and correspond precisely to systems for
which decoupling is valid. This is consistent with the theorem: $F = 0
\Leftrightarrow$ decoupling works.
\end{replybox}

% =================================================================
\section{The Davis Law as Special Case}
\label{sec:davis-law}
% =================================================================

\begin{proposition}[Davis Law encodes the Non-Decoupling Theorem]
\label{prop:davis-law}
The Davis Law $C = \tau / K$ is a scalar projection of
Theorem~\ref{thm:main} onto a one-dimensional summary statistic. Setting
$K = \norm{\Omega}$ (the operator norm of the curvature at a point),
$\tau$ the \emph{tolerance} (regularization capacity of the geometry), and
$C$ the \emph{representational capacity} of the model, the following
correspondences hold:

\begin{center}
\begin{tabular}{lll}
\hline
\textbf{Assumption} & \textbf{Davis Law} & \textbf{Geometric consequence} \\
\hline
Decoupling ($K = 0$) & $C \to \infty$ & Unphysical capacity: model claims to \\
 & & represent all functions \\
Decoupling ($K = 0$) & $\tau \to 0$ & Zero tolerance: catastrophic OOD failure \\
Geometry ($K > 0$) & $C = \tau/K < \infty$ & Finite, regularized capacity \\
Max curvature & $C \to 0$ & Capacity collapses: failure near singularities \\
\hline
\end{tabular}
\end{center}

\end{proposition}

\begin{proof}
The decoupling assumption $\omega = 0$ implies $\Omega = 0$, i.e., $K = 0$.
Substituting into $C = \tau/K$: if $\tau > 0$ (i.e., the model has any
representational tolerance), then $C \to \infty$. This is the 
overconfidence failure mode: the model asserts it can represent every
function. Empirically this corresponds to overfitting in distribution.

Alternatively, maintaining $C$ finite requires $\tau \to 0$ as $K \to 0$,
meaning the model's tolerance for distributional shift collapses to zero:
perfect in-sample, catastrophic out-of-sample. This is the brittleness
failure mode.

When $K > 0$, the geometry provides the regularizer: $C = \tau / K$ is
finite and determined by the ratio of tolerance to curvature. Large
curvature ($K \gg 1$) at a point means the geometry is rapidly changing,
the model's capacity there is small, and errors concentrate---consistent
with Part III of Theorem~\ref{thm:main}.
\end{proof}

\begin{auditbox}
\textbf{Objection 6.} Proposition~\ref{prop:davis-law} sets $K =
\norm{\Omega}$ (a 2-form norm) and identifies this with the scalar $K$ in
$C = \tau/K$. But the Davis Law $C = \tau/K$ was originally stated with $K$
as the sectional curvature of a Riemannian manifold, not the curvature of a
principal bundle connection. These are different geometric objects. The
identification needs justification, or the Davis Law must be re-derived in
the bundle setting.
\end{auditbox}

\begin{replybox}
This is a legitimate gap. The reconciliation runs as follows.

For an associated vector bundle $E = P \times_\rho V$ (where $\rho: G \to
GL(V)$ is a representation), the connection $\omega$ on $P$ induces a
covariant derivative $\nabla$ on $E$. The curvature of $\nabla$ is a
2-form $R^\nabla \in \Omega^2(M, \End(E))$, which contracts to give
sectional curvatures $K(\sigma) = \langle R^\nabla(e_1,e_2)e_2, e_1\rangle$
for 2-planes $\sigma = \mathrm{span}(e_1,e_2)$.

For the specific case of the tangent bundle ($E = TM$, $\rho = $ adjoint),
this recovers the Riemannian sectional curvature. So the Davis Law with $K =
$ sectional curvature is a special case of the bundle framework with $P =
\mathrm{Fr}(TM)$ (the frame bundle).

For general bundles, $K$ in the Davis Law should be interpreted as the norm
of the induced connection curvature on the associated bundle relevant to the
application. For Marcella: $E$ is the token-embedding bundle, $K =
\norm{R^\nabla}$ is the norm of its induced curvature. The Davis Law is then:
$C = \tau / \norm{R^\nabla}$, which is the bundle version stated in
Proposition~\ref{prop:davis-law}.

A full derivation of the Davis Law from the bundle setting is deferred to a
companion paper.
\end{replybox}

% =================================================================
\section{Master Statement}
\label{sec:master}
% =================================================================

\begin{tcolorbox}[colback=goldbg, colframe=goldline, boxrule=1.5pt,
  title={\bfseries Master Statement}]
Under the standing hypotheses (compact oriented simply-connected $M$,
compact connected $G$, system $S$ faithfully modeled as principal
$G$-bundle with connection $\omega$):

\begin{multline*}
  \Omega \not\equiv 0
  \;\implies\;
  \mathcal{E}(S_{\mathrm{flat}}) \geq c\,\mathcal{S}_{\mathrm{YM}}(\omega) > 0
  \;\implies\;
  \Hol(\omega) \neq \{e\} \\
  \implies\;
  \text{flat models are topologically prohibited (if $c_k \neq 0$).}
\end{multline*}

\smallskip
Each implication is strict. The first is quantitative (Lemma~\ref{lem:ym-bound}).
The second is the Ambrose-Singer theorem (Corollary~\ref{cor:holonomy-memory}).
The third requires non-trivial characteristic classes
(Corollary~\ref{cor:topo-prohibit}).

\smallskip
\textbf{The converse chain:}
\[
  \Omega \equiv 0
  \;\iff\;
  \mathcal{E}(S_{\mathrm{flat}}) = 0
  \;\iff\;
  \Hol(\omega) = \{e\}
  \;\iff\;
  \text{decoupling is valid.}
\]
\end{tcolorbox}

% =================================================================
\section{Scope of Application}
\label{sec:scope}
% =================================================================

The theorem applies, \emph{subject to the fiber bundle representation
assumption being empirically supported}, to any system where:

\begin{enumerate}
  \item Components are assumed independent: economics (parallel agents),
    epidemiology (parallel individuals), neuroscience (parallel brain
    regions), flat AI (parallel attention heads, additive residual stream).

  \item State evolution is modeled without coupling: Euclidean attention,
    additive architectures, standard transformer residual stream.

  \item Long-range dependencies are added on top of flat structure rather
    than encoded geometrically: positional encodings, learned biases,
    post-hoc fine-tuning.
\end{enumerate}

The diagnostic procedure is:
\begin{enumerate}[label=\arabic*.]
  \item Identify the Decoupling Assumption: what two things does the model
    treat as independent parallel processes?
  \item Construct or estimate the curvature $\Omega$ of the correct bundle.
  \item Evaluate $\mathcal{S}_{\mathrm{YM}}(\omega)$. This is the size of
    the bill coming due.
  \item For non-trivial $c_k$: the bill cannot be paid with any flat model.
    The solution space requires geometry.
\end{enumerate}

% =================================================================
\section{Open Questions and Future Work}
\label{sec:open}
% =================================================================

\begin{enumerate}
  \item \textbf{Representation theorem.} Prove (or characterize the
    conditions under which) transformer architectures, financial networks,
    and epidemiological contact graphs admit faithful fiber bundle
    descriptions. This is the missing Part (1) identified in Objection~4.

  \item \textbf{Tightness of the Yang-Mills bound.} Is
    $\mathcal{E}(S_{\mathrm{flat}}) = \mathcal{S}_{\mathrm{YM}}(\omega)$
    achievable, or is the bound only a lower bound? A tight bound would give
    the exact holonomy debt.

  \item \textbf{Davis Law derivation.} Derive $C = \tau/K$ rigorously from
    the bundle framework (Proposition~\ref{prop:davis-law} deferred this).

  \item \textbf{Non-simply-connected $M$.} Characterize the failure modes
    when $\pi_1(M) \neq 0$: the interplay between topological ($c_k$) and
    fundamental-group ($\pi_1$) obstructions.

  \item \textbf{Infinite-dimensional bundles.} For Marcella and general
    neural architectures, $M$ is the token-sequence manifold and $P$ is an
    infinite-dimensional bundle. Uhlenbeck's compactness theory extends to
    this setting only under restrictive conditions. A functional-analytic
    treatment is needed.
\end{enumerate}

% =================================================================
\section*{Acknowledgments}
% =================================================================

The fiber bundle interpretation emerged from the empirical failure of V3RG
and the subsequent physics-guided redesign described in the companion
experimental report. The adversarial audit in this paper was conducted by
the author against her own theorem. The author thanks the referee who
will object to Objection~4.

\bibliographystyle{amsalpha}
\begin{thebibliography}{99}

\bibitem{AS53}
W.~Ambrose and I.~M.~Singer,
\textit{A theorem on holonomy},
Trans.\ Amer.\ Math.\ Soc.\ \textbf{75} (1953), 428--443.

\bibitem{Berry84}
M.~V.~Berry,
\textit{Quantal phase factors accompanying adiabatic changes},
Proc.\ R.\ Soc.\ Lond.\ A \textbf{392} (1984), 45--57.

\bibitem{KN63}
S.~Kobayashi and K.~Nomizu,
\textit{Foundations of Differential Geometry}, Vols.~I--II,
Wiley Interscience, 1963.

\bibitem{MS74}
J.~Milnor and J.~Stasheff,
\textit{Characteristic Classes},
Ann.\ of Math.\ Studies \textbf{76}, Princeton Univ.\ Press, 1974.

\bibitem{SW09}
U.~Schreiber and K.~Waldorf,
\textit{Parallel transport and functors},
J.\ Homotopy Relat.\ Struct.\ \textbf{4} (2009), 187--244.

\bibitem{TKNN82}
D.~J.~Thouless, M.~Kohmoto, M.~P.~Nightingale, and M.~den~Nijs,
\textit{Quantized Hall conductance in a two-dimensional periodic potential},
Phys.\ Rev.\ Lett.\ \textbf{49} (1982), 405--408.

\bibitem{Uhl82}
K.~Uhlenbeck,
\textit{Connections with $L^p$ bounds on curvature},
Comm.\ Math.\ Phys.\ \textbf{83} (1982), 31--42.

\bibitem{Weh04}
K.~Wehrheim,
\textit{Uhlenbeck Compactness},
EMS Series of Lectures in Mathematics, 2004.

\end{thebibliography}

\end{document}